
\documentclass[11pt,a4paper]{article}

% Packages
\usepackage[utf8]{inputenc}
\usepackage[T1]{fontenc}
\usepackage{amsmath,amssymb,amsthm}
\usepackage{graphicx}
\usepackage{hyperref}
\usepackage[margin=1in]{geometry}
\usepackage{xfrac}

\theoremstyle{plain}
\newtheorem{theorem}{Theorem}[section]
\theoremstyle{definition}
\newtheorem{definition}[theorem]{Definition}
\newtheorem{remark}{Remark}


% Document information
\title{Document Title}
\author{Author Name}
\date{\today}

\begin{document}

\maketitle

\section{Introduction}

\subsection{Matérn covariance function and its parametrizations}

From \cite{Rasmussen2006Gaussian}, the Matérn covariance function is defined as 
\begin{align}
    r(h) = \dfrac{\sigma^2}{2^{\nu-1}\Gamma(\nu)}\left(\dfrac{\sqrt{2\nu}h}{\rho}\right)^\nu K_\nu\left(\dfrac{\sqrt{2\nu}h}{\rho}\right)
\end{align}
and its corresponding spectral density is 
\begin{align}
\label{eq:spectral_density_2pi}
    \hat{f}(\omega) = \dfrac{\sigma^2\sqrt{4\pi}\Gamma(\nu+\sfrac{1}{2})}{\Gamma(\nu)}\left(\dfrac{2\nu}{\rho^2}\right)^\nu \left(\dfrac{2\nu}{\rho^2}+4\pi^2\omega^2\right)^{-(\nu+\sfrac{1}{2})}
\end{align}
computed by using the Fourier transform pair parametrization 
\begin{align}
    \hat{f}(\omega) = \int_{\mathbb{R}}e^{-i2\pi\omega h}r(h)dh \iff r(h)=\int_{\mathbb{R}}e^{i2\pi\omega h}\hat{f}(\omega)d\omega
\end{align}
But actually we want the pair parametrization
\begin{align}
    \tilde{f}(\omega) = \int_{\mathbb{R}}e^{-i\omega h}r(h)dh \iff r(h)=\dfrac{1}{2\pi}\int_{\mathbb{R}}e^{i\omega h}\tilde{f}(\omega)d\omega
\end{align}
For that, we simply do
\begin{align}
\label{eq:spectral_density}
    \tilde{f}(\omega) = \hat{f}\left(\dfrac{\omega}{2\pi}\right) = \dfrac{\sigma^2\sqrt{4\pi}\Gamma(\nu+\sfrac{1}{2})}{\Gamma(\nu)}\left(\dfrac{2\nu}{\rho^2}\right)^\nu \left(\dfrac{2\nu}{\rho^2}+\omega^2\right)^{-(\nu+\sfrac{1}{2})}
\end{align}
But actually we apparently want this Fourier pair parametrization
\begin{align}
\label{eq:this_is_THE_Fourier_conv_used}
    \overline{f}(\omega) = \mathfrak{F}[r(h)](\omega) =  \dfrac{1}{\sqrt{2\pi}} \int_{\mathbb{R}}e^{-i\omega h}r(h)dh \iff r(h)= \mathfrak{F}^{-1}[\overline{f}(\omega)](h)= \dfrac{1}{\sqrt{2\pi}}\int_{\mathbb{R}}e^{i\omega h}\overline{f}(\omega)d\omega
\end{align}
For that, we simply do
\begin{align}
   \overline{f}(\omega) = \dfrac{1}{\sqrt{2\pi}} \tilde{f}(\omega)  = \dfrac{\sigma^2\sqrt{2}\Gamma(\nu+\sfrac{1}{2})}{\Gamma(\nu)}\left(\dfrac{2\nu}{\rho^2}\right)^\nu \left(\dfrac{2\nu}{\rho^2}+\omega^2\right)^{-(\nu+\sfrac{1}{2})}
\end{align}
Recall the scaling property of the Fourier transform given by 
\begin{align}
    r(ah)\iff \dfrac{1}{|a|}\overline{f}\left(\dfrac{\omega}{a}\right)
\end{align}
Now let $\varrho(h) = r(2h)$ and $\kappa = \sqrt{8\nu}/\rho$, then
\begin{align}
    \varrho(h) & = \dfrac{\sigma^2}{2^{\nu-1}\Gamma(\nu)}\left(\dfrac{\sqrt{8\nu}h}{\rho}\right)^\nu K_\nu\left(\dfrac{\sqrt{8\nu}h}{\rho}\right)\\
    & = \dfrac{\sigma^2}{2^{\nu-1}\Gamma(\nu)}(\kappa h)^\nu K_\nu(\kappa h)
\end{align}
which is the familiar form of the Matérn covariance function.
Just for the record, if we have it in terms of $\tau$, then 
\begin{align}
    \varrho(h) & = \dfrac{1}{2^{\nu-1}\Gamma(\nu)} \dfrac{\Gamma(\nu)}{\Gamma(\nu+\sfrac{1}{2})\sqrt{4\pi}\kappa^{2\nu}\tau^2}(\kappa h)^\nu K_\nu(\kappa h)\\
    & = \dfrac{1}{2^{\nu-1}\Gamma(\nu+\sfrac{1}{2})\sqrt{4\pi}\kappa^{2\nu}\tau^2}(\kappa h)^\nu K_\nu(\kappa h)
\end{align}
Setting $\nu = \tfrac{1}{2}$, we obtain
\begin{align}
\varrho(h)
&=
\dfrac{1}{2^{-1/2}\Gamma(1)\sqrt{4\pi}\kappa^{1}\tau^2}
(\kappa h)^{1/2} K_{1/2}(\kappa h).
\end{align}

Using the identities
\begin{align}
\Gamma(1) &= 1, \\
K_{1/2}(x) &= \sqrt{\frac{\pi}{2x}} e^{-x},
\end{align}
we get
\begin{align}
(\kappa h)^{1/2} K_{1/2}(\kappa h)
&=
(\kappa h)^{1/2}
\sqrt{\frac{\pi}{2\kappa h}}
e^{-\kappa h}
=
\sqrt{\frac{\pi}{2}}\, e^{-\kappa h}.
\end{align}

Substituting back,
\begin{align}
\varrho(h)
&=
\dfrac{1}{2^{-1/2}\sqrt{4\pi}\kappa\tau^2}
\sqrt{\frac{\pi}{2}}\, e^{-\kappa h}.
\end{align}

Now simplify the constants:
\begin{align}
2^{-1/2}\sqrt{4\pi}
&=
\frac{1}{\sqrt{2}} \cdot 2\sqrt{\pi}
=
\sqrt{2\pi},
\end{align}
so that
\begin{align}
\varrho(h)
&=
\frac{1}{\sqrt{2\pi}\,\kappa\tau^2}
\sqrt{\frac{\pi}{2}}\, e^{-\kappa h}
=
\frac{1}{2\kappa \tau^2} e^{-\kappa h}.
\end{align}

By the scaling property, the corresponding spectral density to $\varrho(h)$ is $f(\omega) = \dfrac{1}{2}\overline{f}\left(\dfrac{\omega}{2}\right)$ given by
\begin{align}
    f(\omega) & = \dfrac{\sigma^2\sqrt{2}\Gamma(\nu+\sfrac{1}{2})}{2\Gamma(\nu)}\left(\dfrac{2\nu}{\rho^2}\right)^\nu \left(\dfrac{2\nu}{\rho^2}+\dfrac{\omega^2}{4}\right)^{-(\nu+\sfrac{1}{2})}\\
   & = \dfrac{\sigma^2\sqrt{2}\Gamma(\nu+\sfrac{1}{2})}{2\Gamma(\nu)}\left(\dfrac{2\nu}{\rho^2}\right)^\nu \left(\dfrac{1}{4}\right)^{-(\nu+\sfrac{1}{2})}\left(\dfrac{8\nu}{\rho^2}+\omega^2\right)^{-(\nu+\sfrac{1}{2})}
\end{align}
Since $\left(\dfrac{1}{4}\right)^{-(\nu+\sfrac{1}{2})} = 2\cdot 4^\nu$ and $\kappa^2 = 8\nu/\rho^2$ we have that
\begin{align}
    f(\omega) &= \dfrac{\sigma^2\sqrt{2}\Gamma(\nu+\sfrac{1}{2})}{\Gamma(\nu)}\left(\dfrac{8\nu}{\rho^2}\right)^\nu \left(\dfrac{8\nu}{\rho^2}+\omega^2\right)^{-(\nu+\sfrac{1}{2})}\\
    &= \dfrac{\sigma^2\sqrt{2}\Gamma(\nu+\sfrac{1}{2})\kappa^{2\nu}}{\Gamma(\nu)}\left(\kappa^2+\omega^2\right)^{-(\nu+\sfrac{1}{2})}
\end{align}
Let $A = A_{(\nu,\kappa)} = \dfrac{\sqrt{2}\Gamma(\nu+\sfrac{1}{2})\kappa^{2\nu}}{\Gamma(\nu)}$ so this matches with \cite{Bolin2025Linear}. If $\sigma^2 =  \dfrac{\Gamma(\nu)}{\Gamma(\nu+\sfrac{1}{2})\sqrt{4\pi}\kappa^{2\nu}\tau^2}$, then $A\sigma^2 = (2\pi)^{-\sfrac{1}{2}}\tau^{-2}$. The moral of the story is that 
\begin{align}
    f(\omega) = A\sigma^2\left(\kappa^2+\omega^2\right)^{-(\nu+\sfrac{1}{2})}\iff \varrho(h) = \dfrac{\sigma^2}{2^{\nu-1}\Gamma(\nu)}(\kappa h)^\nu K_\nu(\kappa h), 
\end{align}
Putting in terms of $\alpha$, we instead have
\begin{align}
    f(\omega) = A\sigma^2\left(\kappa^2+\omega^2\right)^{-\alpha}\iff \varrho(h) = \dfrac{\sigma^2}{2^{\alpha-\sfrac{1}{2}-1}\Gamma(\alpha-\sfrac{1}{2})}(\kappa h)^{\alpha-\sfrac{1}{2}} K_{\alpha-\sfrac{1}{2}}(\kappa h), 
\end{align}

If we drop $\sigma^2$, we still have the Fourier pair
\begin{align}
    f_\star(\omega) = A\left(\kappa^2+\omega^2\right)^{-(\nu+\sfrac{1}{2})}\iff \varrho_\star(h) = \dfrac{1}{2^{\nu-1}\Gamma(\nu)}(\kappa h)^\nu K_\nu(\kappa h), 
\end{align}


\newpage
\subsection{Computational quirks}
The \texttt{MetricGraph} package actually codifies
\begin{align}
\overline{\varrho}(h; \kappa, \nu, \texttt{cte})=\dfrac{\texttt{cte}^{2}}{2^{\nu-1} \Gamma(\nu)}(\kappa|h|)^\nu K_\nu(\kappa|h|)
\end{align}
So if you want something like 
\begin{align}
    \dfrac{c_\alpha \sqrt{\pi}}{\sqrt{1-p_i}}\varrho(h; \kappa, \nu, \sigma) = \dfrac{c_\alpha \sqrt{\pi}}{\sqrt{1-p_i}} \dfrac{\sigma^{2}}{2^{\nu-1} \Gamma(\nu)}(\kappa|h|)^\nu K_\nu(\kappa|h|)
\end{align}
computationally you should actually do
\begin{align}
    \sigma^2 \overline{\varrho}\left(h; \kappa, \nu, \sqrt{ \dfrac{c_\alpha \sqrt{\pi}}{\sqrt{1-p_i}}}\right)
\end{align}


\subsection{Preliminaries}

We consider the Matérn covariance function
\begin{align}
\label{eq:Matern_cov_fun}
\varrho(h; \kappa, \nu, \sigma)=\dfrac{\sigma^{2}}{2^{\nu-1} \Gamma(\nu)}(\kappa|h|)^\nu K_\nu(\kappa|h|),\quad \sigma^2 = A^{-1}(2\pi)^{-\sfrac{1}{2}}\tau^{-2},
\end{align}
where $\rho = \sqrt{8\nu}/\kappa$ and $\nu = \alpha-\sfrac{1}{2}$
and the SPDE
\begin{equation}
\label{eq:Matern_spde}
	(\kappa^2 - \Delta_\Gamma)^{\sfrac{\alpha}{2}} (\tau u) = \mathcal{W}, \qquad \text{on $\Gamma$},
\end{equation}


\subsection{Computing spectral density from SPDE}

We take the Fourier transform \eqref{eq:this_is_THE_Fourier_conv_used} to \eqref{eq:Matern_spde} to get
\begin{align}
    (\kappa^2+\omega^2)^{\sfrac{\alpha}{2}}\tau\hat{u}(\omega) = \widehat{\mathcal{W}}(\omega)\implies
    f(\omega) = \dfrac{1}{\tau^2\sqrt{2\pi}}(\kappa^2+\omega^2)^{-\alpha}
\end{align}

Let $L = (\kappa^2 - \Delta_\Gamma)$. The solution $u$ is a centered Gaussian random field with covariance operator $\mathcal{C} = \tau^{-2}L^{-\alpha}$. We consider the approximation
\begin{align}
    \mathcal{C}_\alpha = \tau^{-2}L^{-\lfloor\alpha\rfloor}P_m(L)Q_m(L)^{-1}
\end{align}
where 
\begin{align}
    x^{\{\alpha\}}\approx \dfrac{p_m(x)}{q_m(x)} = \dfrac{\sum_{i=0}^m a_ix^{i}}{\sum_{j=0}^m b_jx^{j}} = \dfrac{P_m(x^{-1})}{Q_m(x^{-1})} =r_m(x)
\end{align}
or
\begin{align}
\label{eq:rat_approx}
    x^{-\{\alpha\}}\approx \dfrac{p_m(x^{-1})}{q_m(x^{-1})} = \dfrac{\sum_{i=0}^m a_ix^{m-i}}{\sum_{j=0}^m b_jx^{m-j}} = \dfrac{P_m(x)}{Q_m(x)} = R_m(x)
\end{align}
We further have that
\begin{align}
\label{eq:partial_frac_dec}
    \dfrac{P_m(x)}{Q_m(x)} = \dfrac{p_m(x^{-1})}{q_m(x^{-1})} =k + \sum_{i=1}^mr_i(x-p_i)^{-1}.
\end{align}
Therefore
\begin{align}
    \mathcal{C}_\alpha = \tau^{-2}L^{-\lfloor\alpha\rfloor}\left(kI + \sum_{i=1}^mr_i(L-p_iI)^{-1}\right) = \tau^{-2}kL^{-\lfloor\alpha\rfloor} + \sum_{i=1}^m\tau^{-2}r_iL^{-\lfloor\alpha\rfloor}(L-p_iI)^{-1}
\end{align}
Since $k,r_i>0$ and $p_i<0$, $\mathcal{C}_0 = \tau^{-2}kL^{-\lfloor\alpha\rfloor}$ and $\mathcal{C}_i = \tau^{-2}r_iL^{-\lfloor\alpha\rfloor}(L-p_iI)^{-1}$ are valid covariance operators so that $\mathcal{C}_\alpha = \mathcal{C}_0+\sum_{i=1}^m\mathcal{C}_i$. Then the approximation $u_\alpha$ has decomposition
\begin{align}
    u_\alpha = u_0+\sum_{i=1}^mu_i
\end{align}
where $u_0\sim N(0,\mathcal{C}_0)$ and $u_i\sim N(0,\mathcal{C}_i)$. Moreover, they satisfy
\begin{align}
\label{eq:WM_process}
   \mathcal{C}_0^{-\sfrac{1}{2}}u_0= L^{\sfrac{\lfloor\alpha\rfloor}{2}}(\tau k^{-\sfrac{1}{2}}u_0) = \mathcal{W}, \qquad \text{on $\Gamma$},
\end{align}
and
\begin{align}
\label{eq:shifted_WM_process}
   \mathcal{C}_i^{-\sfrac{1}{2}}u_i =  (L^{\lfloor\alpha\rfloor}(L-p_iI))^{\sfrac{1}{2}}(\tau r_i^{-\sfrac{1}{2}}u_i) = \mathcal{W}, \qquad \text{on $\Gamma$}
\end{align}
\begin{definition}
    $u_i$ in \eqref{eq:shifted_WM_process} is called a shifted Whittle-Matérn process on $\Gamma$.
\end{definition}
\begin{remark}
    Observe that \eqref{eq:WM_process} are Whittle-Matérn process where the power $\lfloor\alpha\rfloor$ is integer.
\end{remark}
Observe now that \eqref{eq:Matern_cov_fun} has spectral density 
\begin{align}
f(\omega) = A\sigma^2(\kappa^2+\omega^2)^{-(\nu+\sfrac{1}{2})},\quad A = \sigma^{-2}(2\pi)^{-\sfrac{1}{2}}\tau^{-2}
\end{align}
which can be also written as 
\begin{align}
    f(\omega) = A\sigma^2\kappa^{-2\alpha}(1+\kappa^{-2}\omega^2)^{-\alpha} = A\sigma^2\kappa^{-2\alpha}(1+\kappa^{-2}\omega^2)^{-\lfloor\alpha\rfloor} (1+\kappa^{-2}\omega^2)^{-\{\alpha\}}
\end{align}
By \eqref{eq:rat_approx}, we can approximate $f$ by $f_\alpha$ given by
\begin{align}
    f_\alpha(\omega) = A\sigma^2\kappa^{-2\alpha}\dfrac{1}{(1+\kappa^{-2}\omega^2)^{\lfloor\alpha\rfloor} }\dfrac{p_m(1+\kappa^{-2}\omega^2)}{q_m(1+\kappa^{-2}\omega^2)}
\end{align}
By \eqref{eq:partial_frac_dec}, we have that 
\begin{align}
    f_\alpha(\omega)& = A\sigma^2\kappa^{-2\alpha}\dfrac{1}{(1+\kappa^{-2}\omega^2)^{\lfloor\alpha\rfloor} }\left(k + \sum_{i=1}^mr_i(1+\kappa^{-2}\omega^2-p_i)^{-1}\right)\\
    & = \dfrac{A\sigma^2\kappa^{-2\alpha}k}{(1+\kappa^{-2}\omega^2)^{\lfloor\alpha\rfloor}}+ \sum_{i=1}^m\dfrac{A\sigma^2\kappa^{-2\alpha}r_i}{(1+\kappa^{-2}\omega^2)^{\lfloor\alpha\rfloor}(1+\kappa^{-2}\omega^2-p_i)}
\end{align}
Let 
\begin{align}
    f_0(\omega) = \dfrac{A\sigma^2\kappa^{-2\alpha}k}{(1+\kappa^{-2}\omega^2)^{\lfloor\alpha\rfloor}}\quad\text{ and }\quad f_i(\omega) = \dfrac{A\sigma^2\kappa^{-2\alpha}r_i}{(1+\kappa^{-2}\omega^2)^{\lfloor\alpha\rfloor}(1+\kappa^{-2}\omega^2-p_i)}
\end{align}
then we can write
\begin{align}
    f_\alpha(\omega) = f_0(\omega) + \sum_{i=1}^mf_i(\omega),
\end{align}
which is a valid decomposition of spectral densities.
We now can equivalently write 
\begin{align}
    f_0(\omega) = \dfrac{A\sigma^2\kappa^{-2\alpha}k\kappa^{2\lfloor\alpha\rfloor}}{(\kappa^2+\omega^2)^{\lfloor\alpha\rfloor}}\quad \text{ and }\quad f_i(\omega) = \dfrac{A\sigma^2\kappa^{-2\alpha}r_i\kappa^{2\lfloor\alpha\rfloor+2}}{(\kappa^{2}+\omega^2)^{\lfloor\alpha\rfloor}(\kappa^{2}(1-p_i)+\omega^2)}
\end{align}
Now for $k\in\mathbb{N}$ and $p\in\mathbb{R}$, we have that
\begin{align}
    \dfrac{1}{(\kappa^{2}+\omega^2)^{k}(\kappa^{2}(1-p)+\omega^2)}& = \dfrac{1}{(\kappa^{2}p)^{k}(\kappa^{2}(1-p)+\omega^2)} - \sum_{j=1}^k  \dfrac{1}{(\kappa^{2}p)^{k+1-j}(\kappa^{2}+\omega^2)^j}
\end{align}
Indeed, let $r = \dfrac{\kappa^2p}{\kappa^{2}+\omega^2}$, then $1-r = \dfrac{\kappa^2(1-p)+\omega^2}{\kappa^2+\omega^2}$. Hence
\begin{align*}
    \sum_{j=1}^k  \dfrac{1}{(\kappa^{2}p)^{k+1-j}(\kappa^{2}+\omega^2)^j} & = \dfrac{1}{(\kappa^{2}p)^{k+1}} \sum_{j=1}^k  \left(\dfrac{\kappa^2p}{\kappa^{2}+\omega^2}\right)^j\\
    & = \dfrac{1}{(\kappa^{2}p)^{k+1}} \dfrac{r(1-r^k)}{1-r}\\
    & = \dfrac{1}{(\kappa^{2}p)^{k+1}} \dfrac{\dfrac{\kappa^2p}{\kappa^{2}+\omega^2}\left(1-\left(\dfrac{\kappa^2p}{\kappa^{2}+\omega^2}\right)^k\right)}{\dfrac{\kappa^2(1-p)+\omega^2}{\kappa^2+\omega^2}}\\
    & = \dfrac{1-\left(\dfrac{\kappa^2p}{\kappa^{2}+\omega^2}\right)^k}{(\kappa^{2}p)^{k}(\kappa^2(1-p)+\omega^2)}\\
    & = \dfrac{1}{(\kappa^{2}p)^{k}(\kappa^2(1-p)+\omega^2)} - \dfrac{1}{(\kappa^2+\omega^2)^k(\kappa^2(1-p)+\omega^2)}
\end{align*}
Therefore
\begin{align}
    f_i(\omega) &= \dfrac{A\sigma^2\kappa^{-2\alpha}r_i\kappa^{2\lfloor\alpha\rfloor+2}}{(\kappa^{2}p_i)^{\lfloor\alpha\rfloor}(\kappa^{2}(1-p_i)+\omega^2)}\\
    &- \sum_{j=1}^{\lfloor\alpha\rfloor} \dfrac{A\sigma^2\kappa^{-2\alpha}r_i\kappa^{2\lfloor\alpha\rfloor+2}}{(\kappa^{2}p_i)^{\lfloor\alpha\rfloor+1-j}(\kappa^{2}+\omega^2)^j}\\
    &= \dfrac{A_{(\nu,\kappa)}\kappa^{-2\alpha}r_i\kappa^{2\lfloor\alpha\rfloor+2}}{(\kappa^{2}p_i)^{\lfloor\alpha\rfloor}}\sigma^2(\kappa^{2}(1-p_i)+\omega^2)^{-1}\\
    &- \sum_{j=1}^{\lfloor\alpha\rfloor} \dfrac{A_{(\nu,\kappa)}\kappa^{-2\alpha}r_i\kappa^{2\lfloor\alpha\rfloor+2}}{(\kappa^{2}p_i)^{\lfloor\alpha\rfloor+1-j}}\sigma^2(\kappa^{2}+\omega^2)^{-j}\\
    %
    &= \dfrac{A_{(\nu,\kappa)}\kappa^{-2\alpha}r_i\kappa^{2\lfloor\alpha\rfloor+2}}{(\kappa^{2}p_i)^{\lfloor\alpha\rfloor}A_{(\sfrac{1}{2},\kappa\sqrt{1-p_i})}}A_{(\sfrac{1}{2},\kappa\sqrt{1-p_i})}\sigma^2(\kappa^{2}(1-p_i)+\omega^2)^{-(\sfrac{1}{2}+\sfrac{1}{2})}\\
    &- \sum_{j=1}^{\lfloor\alpha\rfloor} \dfrac{A_{(\nu,\kappa)}\kappa^{-2\alpha}r_i\kappa^{2\lfloor\alpha\rfloor+2}}{(\kappa^{2}p_i)^{\lfloor\alpha\rfloor+1-j}A_{(j-\sfrac{1}{2},\kappa)}}A_{(j-\sfrac{1}{2},\kappa)}\sigma^2(\kappa^{2}+\omega^2)^{-(j-\sfrac{1}{2}+\sfrac{1}{2})}\\
    %
    & = f_{i0}(\omega)- \sum_{j=1}^{\lfloor\alpha\rfloor} f_{ij}(\omega)
\end{align}

\begin{align}
\begin{cases}
    \varrho(h) &= \mathfrak{F}^{-1}[f(\omega)](h)= \dfrac{\sigma^2}{2^{\nu-1}\Gamma(\nu)}(\kappa h)^\nu K_\nu(\kappa h) \\
    f(\omega)& = A_{(\nu,\kappa)}\sigma^2\left(\kappa^2+\omega^2\right)^{-(\nu+\sfrac{1}{2})}\\
    A_{(\nu,\kappa)} & = \dfrac{\sqrt{2}\Gamma(\nu+\sfrac{1}{2})\kappa^{2\nu}}{\Gamma(\nu)} = \dfrac{\sqrt{2}\Gamma(\alpha)\kappa^{2\nu}}{\Gamma(\alpha-\sfrac{1}{2})}\\
    \sigma^2 &= A_{(\nu,\kappa)}^{-1}(2\pi)^{-\sfrac{1}{2}}\tau^{-2}\\
    \rho & = \sqrt{8\nu}/\kappa\\
    \nu &= \alpha-\sfrac{1}{2}
\end{cases}
\end{align}

Similarly for $f_0$
\begin{align}
    f_0(\omega) = A_{(\nu,\kappa)}\kappa^{-2\alpha}k\kappa^{2\lfloor\alpha\rfloor}\sigma^2(\kappa^2+\omega^2)^{-\lfloor\alpha\rfloor} = \dfrac{A_{(\nu,\kappa)}\kappa^{-2\alpha}k\kappa^{2\lfloor\alpha\rfloor}}{A_{(\lfloor\alpha\rfloor-\sfrac{1}{2},\kappa)}}A_{(\lfloor\alpha\rfloor-\sfrac{1}{2},\kappa)}\sigma^2(\kappa^2+\omega^2)^{-(\lfloor\alpha\rfloor-\sfrac{1}{2}+\sfrac{1}{2})}
\end{align}
By taking the inverse Fourier transform, we arrive at
\begin{align}
    \varrho_0(h) = \dfrac{A_{(\nu,\kappa)}\kappa^{-2\alpha}k\kappa^{2\lfloor\alpha\rfloor}}{A_{(\lfloor\alpha\rfloor-\sfrac{1}{2},\kappa)}}\varrho(h;\kappa, \lfloor\alpha\rfloor-\sfrac{1}{2}, \sigma)
\end{align}
and
\begin{align}
    \varrho_i(h)  &= \dfrac{A_{(\nu,\kappa)}\kappa^{-2\alpha}r_i\kappa^{2\lfloor\alpha\rfloor+2}}{(\kappa^{2}p_i)^{\lfloor\alpha\rfloor}A_{(\sfrac{1}{2},\kappa\sqrt{1-p_i})}}\varrho(h;\kappa\sqrt{1-p_i}, \sfrac{1}{2}, \sigma)- \sum_{j=1}^{\lfloor\alpha\rfloor} \dfrac{A_{(\nu,\kappa)}\kappa^{-2\alpha}r_i\kappa^{2\lfloor\alpha\rfloor+2}}{(\kappa^{2}p_i)^{\lfloor\alpha\rfloor+1-j}A_{(j-\sfrac{1}{2},\kappa)}}\varrho(h;\kappa, j-\sfrac{1}{2}, \sigma)
\end{align}
Let us simplify the expressions above. We first simplify $\varrho_0$.
\begin{align*}
    \varrho_0(h) &= k\sigma^2\kappa^{-2\alpha}\kappa^{2\lfloor\alpha\rfloor}\dfrac{A_{(\nu,\kappa)}}{A_{(\lfloor\alpha\rfloor-\sfrac{1}{2},\kappa)}}\varrho(h;\kappa, \lfloor\alpha\rfloor-\sfrac{1}{2}, 1)\\
    &= k\sigma^2\kappa^{-2\alpha}\kappa^{2\lfloor\alpha\rfloor}\dfrac{\dfrac{\sqrt{2}\Gamma(\nu+\sfrac{1}{2})\kappa^{2\nu}}{\Gamma(\nu)}}{\dfrac{\sqrt{2}\Gamma(\lfloor\alpha\rfloor-\sfrac{1}{2}+\sfrac{1}{2})\kappa^{2(\lfloor\alpha\rfloor-\sfrac{1}{2})}}{\Gamma(\lfloor\alpha\rfloor-\sfrac{1}{2})}}\varrho(h;\kappa, \lfloor\alpha\rfloor-\sfrac{1}{2}, 1)\\
    %
    &= k\sigma^2\kappa^{-2\alpha}\kappa^{2\lfloor\alpha\rfloor}\kappa^{2\nu}\kappa^{-2(\lfloor\alpha\rfloor-\sfrac{1}{2})}\dfrac{\dfrac{\Gamma(\alpha)}{\Gamma(\alpha-\sfrac{1}{2})}}{\dfrac{\Gamma(\lfloor\alpha\rfloor)}{\Gamma(\lfloor\alpha\rfloor-\sfrac{1}{2})}}\varrho(h;\kappa, \lfloor\alpha\rfloor-\sfrac{1}{2}, 1)\\
    %
    &= k\sigma^2\kappa^{2(-\alpha+\lfloor\alpha\rfloor+\nu-\lfloor\alpha\rfloor+\sfrac{1}{2})}\dfrac{\dfrac{\Gamma(\alpha)}{\Gamma(\alpha-\sfrac{1}{2})}}{\dfrac{\Gamma(\lfloor\alpha\rfloor)}{\Gamma(\lfloor\alpha\rfloor-\sfrac{1}{2})}}\varrho(h;\kappa, \lfloor\alpha\rfloor-\sfrac{1}{2}, 1)\\
    %
    &= k\sigma^2\dfrac{c_\alpha}{c_{\lfloor\alpha\rfloor}}\varrho(h;\kappa, \lfloor\alpha\rfloor-\sfrac{1}{2}, 1)\\
    %
    &= k\sigma^2\varrho\left(h;\kappa, \lfloor\alpha\rfloor-\sfrac{1}{2}, \sqrt{\sfrac{c_\alpha}{c_{\lfloor\alpha\rfloor}}}\right)
\end{align*}
where $c_a = \sfrac{\Gamma(a)}{\Gamma(a-\sfrac{1}{2})}$. 
For simplicity, let
\begin{align}
    \varrho_{i0}(h) = \dfrac{A_{(\nu,\kappa)}\kappa^{-2\alpha}r_i\kappa^{2\lfloor\alpha\rfloor+2}}{(\kappa^{2}p_i)^{\lfloor\alpha\rfloor}A_{(\sfrac{1}{2},\kappa\sqrt{1-p_i})}}\varrho(h;\kappa\sqrt{1-p_i}, \sfrac{1}{2}, \sigma)
\end{align}
and 
\begin{align}
    \varrho_{ij}(h) = \dfrac{A_{(\nu,\kappa)}\kappa^{-2\alpha}r_i\kappa^{2\lfloor\alpha\rfloor+2}}{(\kappa^{2}p_i)^{\lfloor\alpha\rfloor+1-j}A_{(j-\sfrac{1}{2},\kappa)}}\varrho(h;\kappa, j-\sfrac{1}{2}, \sigma)
\end{align}
Let us now simplify $\varrho_{i0}$.
\begin{align*}
    \varrho_{i0}(h) & = r_i\sigma^2\dfrac{1}{p_i^{\lfloor\alpha\rfloor}} \kappa^{-2\alpha}\kappa^{2\lfloor\alpha\rfloor+2}\kappa^{-2\lfloor\alpha\rfloor}\dfrac{A_{(\nu,\kappa)}}{A_{(\sfrac{1}{2},\kappa\sqrt{1-p_i})}}\varrho(h;\kappa\sqrt{1-p_i}, \sfrac{1}{2}, 1)\\
    %
    & = r_i\sigma^2\dfrac{1}{p_i^{\lfloor\alpha\rfloor}} \kappa^{-2\alpha}\kappa^{2\lfloor\alpha\rfloor+2}\kappa^{-2\lfloor\alpha\rfloor}\dfrac{\dfrac{\sqrt{2}\Gamma(\alpha)\kappa^{2\nu}}{\Gamma(\alpha-\sfrac{1}{2})}}{\dfrac{\sqrt{2}\Gamma(\sfrac{1}{2}+\sfrac{1}{2})(\kappa\sqrt{1-p_i})^{2\cdot\sfrac{1}{2}}}{\Gamma(\sfrac{1}{2})}}\varrho(h;\kappa\sqrt{1-p_i}, \sfrac{1}{2}, 1)\\
    %
    & = r_i\sigma^2\dfrac{1}{p_i^{\lfloor\alpha\rfloor}} \kappa^{2(-\alpha+\lfloor\alpha\rfloor+1-\lfloor\alpha\rfloor+\nu-\sfrac{1}{2})}\dfrac{\dfrac{\Gamma(\alpha)}{\Gamma(\alpha-\sfrac{1}{2})}}{\dfrac{\Gamma(\sfrac{1}{2}+\sfrac{1}{2})\sqrt{1-p_i}}{\Gamma(\sfrac{1}{2})}}\varrho(h;\kappa\sqrt{1-p_i}, \sfrac{1}{2}, 1)\\
    %
    & = r_i\sigma^2\dfrac{1}{p_i^{\lfloor\alpha\rfloor}} \dfrac{c_\alpha\sqrt{\pi}}{\sqrt{1-p_i}}\varrho(h;\kappa\sqrt{1-p_i}, \sfrac{1}{2}, 1)\\
    %
    & = r_i\sigma^2\dfrac{1}{p_i^{\lfloor\alpha\rfloor}} \varrho\left(h;\kappa\sqrt{1-p_i}, \sfrac{1}{2}, \sqrt{\sfrac{c_\alpha\sqrt{\pi}}{\sqrt{1-p_i}}}\right)
\end{align*}
since $\Gamma(1)=1$ and $\Gamma(\sfrac{1}{2}) = \sqrt{\pi}$. Now we simplify $\varrho_{ij}$.
\begin{align*}
    \varrho_{ij}(h) & = r_i\sigma^2\kappa^{2(-\alpha+\lfloor\alpha\rfloor+1-\lfloor\alpha\rfloor-1+j)}\dfrac{1}{p_i^{\lfloor\alpha\rfloor+1-j}}\dfrac{A_{(\nu,\kappa)}}{A_{(j-\sfrac{1}{2},\kappa)}}\varrho(h;\kappa, j-\sfrac{1}{2}, 1)\\
    %
    & = r_i\sigma^2\kappa^{2(-\alpha+\lfloor\alpha\rfloor+1-\lfloor\alpha\rfloor-1+j)}\dfrac{1}{p_i^{\lfloor\alpha\rfloor+1-j}}\dfrac{\dfrac{\sqrt{2}\Gamma(\alpha)\kappa^{2\nu}}{\Gamma(\alpha-\sfrac{1}{2})}}{\dfrac{\sqrt{2}\Gamma(j)\kappa^{2(j-\sfrac{1}{2})}}{\Gamma(j-\sfrac{1}{2})}}\varrho(h;\kappa, j-\sfrac{1}{2}, 1)\\
    %
    & = r_i\sigma^2\kappa^{2(-\alpha+\lfloor\alpha\rfloor+1-\lfloor\alpha\rfloor-1+j+\nu-j+\sfrac{1}{2})}\dfrac{1}{p_i^{\lfloor\alpha\rfloor+1-j}}\dfrac{c_\alpha}{c_j}\varrho(h;\kappa, j-\sfrac{1}{2}, 1)\\
    %
    & = r_i\sigma^2\dfrac{1}{p_i^{\lfloor\alpha\rfloor+1-j}}\varrho\left(h;\kappa, j-\sfrac{1}{2}, \sqrt{\sfrac{c_\alpha}{c_j}}\right)
\end{align*}
To sum up, we have arrived at
\begin{align}
   \varrho_\alpha(h) = \varrho_0(h) + \sum_{i=1}^m  \varrho_i(h)
\end{align}
where 
\begin{align}
   \varrho_0(h) =  k\sigma^2\varrho\left(h;\kappa, \lfloor\alpha\rfloor-\sfrac{1}{2}, \sqrt{\sfrac{c_\alpha}{c_{\lfloor\alpha\rfloor}}}\right)
\end{align}
and 
\begin{align}
    \varrho_i(h) = r_i\sigma^2\dfrac{1}{p_i^{\lfloor\alpha\rfloor}} \varrho\left(h;\kappa\sqrt{1-p_i}, \sfrac{1}{2}, \sqrt{\sfrac{c_\alpha\sqrt{\pi}}{\sqrt{1-p_i}}}\right) -  \sum_{j=1}^{\lfloor\alpha\rfloor} r_i\sigma^2\dfrac{1}{p_i^{\lfloor\alpha\rfloor+1-j}}\varrho\left(h;\kappa, j-\sfrac{1}{2}, \sqrt{\sfrac{c_\alpha}{c_j}}\right)
\end{align}

\section{Comparison to rat int paper}

The calculation above is a detailed derivation of the partial-fraction
identity and inverse Fourier transforms underlying Proposition~5 of
\cite{Bolin2025Linear}, equation~(9) on page~7 of
\texttt{rat\_int\_paper.pdf}; its proof appears on pages~28--29.
The final component formulas agree after translating the notation, but
some conventions, qualifications, and inconsistencies must be separated.
The comparison below refers to the supplied PDF, including its printed
formulas, rather than to a potentially revised version of the paper.

\subsection{Notation and agreement of the component formulas}

The component covariance formulas coincide with Proposition~5 after
identifying our residues $r_i$ with the paper's $c_i$ and translating
the Mat\'ern parametrization:
\begin{equation}
\varrho(h;\kappa,\nu,s)
=\varrho^{\mathrm{paper}}(h;\nu,\kappa,s^2).
\end{equation}
Our last argument is a standard deviation, whereas the paper's is a
variance; this explains the square roots in our final arguments.
For later reference, let $q=\lfloor\alpha\rfloor$,
$\kappa_i=\kappa\sqrt{1-p_i}$, and
$M_{\nu,\lambda}(h)=\varrho(h;\lambda,\nu,1)$. For $q\geq1$, the
common formulas are
\begin{align}
\varrho_0(h)
&=k\sigma^2\frac{c_\alpha}{c_q}M_{q-1/2,\kappa}(h),
\label{eq:comparison_rho0}\\
\varrho_i(h)
&=r_i\sigma^2c_\alpha\left[
 \frac{\sqrt{\pi}}{p_i^q\sqrt{1-p_i}}M_{1/2,\kappa_i}(h)
 -\sum_{j=1}^q\frac{M_{j-1/2,\kappa}(h)}{p_i^{q+1-j}c_j}
 \right].
\label{eq:comparison_rhoi}
\end{align}
\subsection{Fourier normalization: an inconsistency in the supplied paper}

Equation~(4) in the paper is wrong: its factor $1/(2\pi)$ must be
$1/\sqrt{2\pi}$ to agree with the printed
$A=\sqrt{2}\,c_\alpha\kappa^{2\alpha-1}$ and Proposition~5.
The corrected formula is
\begin{equation}
r_{m,\alpha}(h)=\frac{1}{\sqrt{2\pi}}
\int_{\mathbb R}e^{i\omega h}f_{m,\alpha}(\omega)\,d\omega.
\end{equation}
This is exactly the convention used in
\eqref{eq:this_is_THE_Fourier_conv_used}; after this correction, there
is no Fourier normalization difference between the two derivations.

\subsection{The range \texorpdfstring{$1/2<\alpha<1$}{1/2 < alpha < 1} and the white-noise component}

For $1/2<\alpha<1$, we have $q=\lfloor\alpha\rfloor=0$. Hence the
zeroth term of the spectral decomposition becomes
\begin{equation}
f_0(\omega)
=\frac{A\sigma^2\kappa^{-2\alpha}k}
 {(1+\kappa^{-2}\omega^2)^{\lfloor\alpha\rfloor}}
=A\sigma^2\kappa^{-2\alpha}k.
\end{equation}
Under the symmetric Fourier convention,
$\mathfrak F^{-1}[1](h)=\sqrt{2\pi}\,\delta_0(h)$, since
$\int_{\mathbb R}e^{i\omega h}\,d\omega=2\pi\delta_0(h)$.
Therefore,
\begin{align}
\varrho_0(h)
&=\frac{1}{\sqrt{2\pi}}\int_{\mathbb R}
 e^{i\omega h}A\sigma^2\kappa^{-2\alpha}k\,d\omega\\
&=\sqrt{2\pi}A\sigma^2\kappa^{-2\alpha}k\,\delta_0(h).
\end{align}
Using $A=\sqrt{2}\,c_\alpha\kappa^{2\alpha-1}$ and
$A\sigma^2=(2\pi)^{-1/2}\tau^{-2}$ gives
\begin{equation}
\varrho_0(h)=b\,\delta_0(h),\qquad
b=\sqrt{2\pi}A\sigma^2\kappa^{-2\alpha}k
 =\frac{k\sigma^2c_\alpha\sqrt{4\pi}}{\kappa}
 =\frac{k}{\tau^2\kappa^{2\alpha}}.
\end{equation}

\subsection{Operator definition and scaling in the present derivation}

There is an inconsistency in the operator paragraph preceding
\eqref{eq:rat_approx}: defining
$L=(\kappa^2-\Delta_\Gamma)^{\alpha/2}$ and then writing
$\mathcal C=\tau^{-2}L^{-\alpha}$ applies the exponent twice. Instead
define the positive self-adjoint base operator, with its graph domain
and vertex conditions specified, by
\begin{equation}
L_0=\kappa^2-\Delta_\Gamma,\qquad
\mathcal C=\tau^{-2}L_0^{-\alpha}.
\end{equation}
Equivalently, retaining the original powered $L$ would require
$\mathcal C=\tau^{-2}L^{-2}$. The paper's Section~5 uses an unpowered
base operator in $L^\beta(\tau u)=\mathcal W$; the correspondence here
is $L=L_0$ and $\beta=\alpha/2$. The covariance of $u$ includes
$\tau^{-2}$ (the paper's displayed $L^{-2\beta}$ suppresses this scale).

A second distinction is that equation~(8) and our spectral calculation
apply the rational approximation to the dimensionless variable
$1+\kappa^{-2}\omega^2$. To use the same coefficients $k,r_i,p_i$ at
operator level, define
\begin{equation}
B=\kappa^{-2}L_0=I-\kappa^{-2}\Delta_\Gamma,\qquad
R_m(x)=k+\sum_{i=1}^m\frac{r_i}{x-p_i}.
\end{equation}
Then the consistent operator approximation and its components are
\begin{align}
\mathcal C_m
&=\tau^{-2}\kappa^{-2\alpha}B^{-q}R_m(B)
 =\mathcal C_0+\sum_{i=1}^m\mathcal C_i,
\label{eq:comparison_cov_operator}\\
\mathcal C_0
&=\tau^{-2}k\kappa^{2q-2\alpha}L_0^{-q},\\
\mathcal C_i
&=\tau^{-2}r_i\kappa^{2q+2-2\alpha}
 L_0^{-q}(L_0-\kappa^2p_iI)^{-1}.
\end{align}
In particular, the shifted operator is
$L_0-\kappa^2p_iI=\kappa_i^2-\Delta_\Gamma$, consistent with the
$\kappa_i$ in Proposition~5. The unscaled operator expressions above
\eqref{eq:WM_process} do not use these same dimensionless coefficients
unless $\kappa=1$ or the coefficients have been rescaled. In the
unscaled variable the equivalent coefficients are
\begin{equation}
\widetilde k=k\kappa^{2q-2\alpha},\qquad
\widetilde r_i=r_i\kappa^{2q+2-2\alpha},\qquad
\widetilde p_i=\kappa^2p_i.
\end{equation}
For example, the component SPDEs consistent with
\eqref{eq:comparison_cov_operator} are
\begin{align}
L_0^{q/2}\bigl(\tau k^{-1/2}\kappa^{\alpha-q}u_0\bigr)
&=\mathcal W_0,\\
\bigl[L_0^q(L_0-\kappa^2p_iI)\bigr]^{1/2}
\bigl(\tau r_i^{-1/2}\kappa^{\alpha-q-1}u_i\bigr)
&=\mathcal W_i.
\end{align}
These scaling issues occur in the earlier operator formulation; the
later spectral derivation already contains the required powers of
$\kappa$ and leads to the matching formulas
\eqref{eq:comparison_rho0}--\eqref{eq:comparison_rhoi}.

\subsection{Stationary interval covariances versus metric-graph covariances}

Proposition~5 concerns stationary processes on $\mathbb R$ restricted
to an interval, so their covariances depend on the lag $h=s-t$.
An ordinary Fourier transform diagonalizes $\kappa^2-\Delta$ on
$\mathbb R$; it does not in general diagonalize a metric-graph
Laplacian with vertex conditions. Consequently, the Fourier step
applied to \eqref{eq:Matern_spde} must be understood as the
whole-line calculation, not as a proof that a graph covariance has
the same stationary lag formula.

For a compact metric graph with orthonormal eigenfunctions $\phi_n$
of $-\Delta_\Gamma$ and eigenvalues $\lambda_n\geq0$, the corresponding covariance
is instead
\begin{equation}
C_\Gamma(s,t)=\tau^{-2}\sum_n
(\kappa^2+\lambda_n)^{-\alpha}\phi_n(s)\phi_n(t),
\end{equation}
with the rational approximation obtained by replacing each spectral
multiplier according to \eqref{eq:comparison_cov_operator}.
This generally depends on both $s$ and $t$ and includes topology and
vertex conditions. Even on an interval, imposing Neumann conditions
defines a different model from restricting a stationary whole-line
Mat\'ern process. The paper discusses operator extensions in
Section~5, but Proposition~5 alone does not establish the graph
construction or the required endpoint precision corrections and
conditioning. Those are additional steps beyond the inverse Fourier
calculation reproduced here.

\subsection{Typographical issues in the paper's proof}

The geometric-sum identity used here is the same one used on page~28.
It requires $p\ne0$; the earlier wording ``$p\in\mathbb R$'' should
include that qualification, which is automatic for the actual poles
$p_i<0$. Using $q$ for the integer order also avoids confusing the
identity's integer $k$ with the partial-fraction constant $k$.

In the final displayed spectral decomposition on page~28 of the paper,
the first resolvent term is printed inside the sum over $j$, and some
denominators use $p$ rather than $p_i$. Read literally, that display
does not follow from equation~(8). The consistent expression is
\begin{align}
f_{m,\alpha}(\omega)
&=A\sigma^2\kappa^{-2\alpha}\left[
 \frac{k\kappa^{2q}}{(\kappa^2+\omega^2)^q}
 +\sum_{i=1}^m r_i\kappa^{2q+2}H_i(\omega)\right],\\
H_i(\omega)
&=\frac{1}{(\kappa^2p_i)^q(\kappa_i^2+\omega^2)}
 -\sum_{j=1}^q
 \frac{1}{(\kappa^2p_i)^{q+1-j}(\kappa^2+\omega^2)^j}.
\end{align}
Thus the first term occurs once for each $i$, outside the inner
$j$-sum. The detailed derivation above already uses this placement
and the indexed poles correctly. The discrepancy is in that printed
proof display, not in the final component formulas of Proposition~5.

\section{Implementation}

\subsection{Case $\alpha\in(0,1)$}

\begin{align}
    \varrho_0(h) = k\sigma^2\dfrac{c_\alpha\sqrt{4\pi}}{\kappa}1_{[h=0]}\implies \mathbf{Q}_0 = \mathbf{I}\cdot \dfrac{1}{(k\sigma^2c_\alpha\sqrt{4\pi})/\kappa}
\end{align}

\begin{align}
    \varrho_i(h) = r_i \dfrac{c_\alpha\sqrt{\pi}}{\sqrt{1-p_i}}\varrho(h;\kappa\sqrt{1-p_i}, \sfrac{1}{2}, \sigma) \implies \mathbf{Q}_i = \dfrac{1}{(r_i c_\alpha\sqrt{\pi})/\sqrt{1-p_i}}\begin{cases}
        \overline{\nu} = \sfrac{1}{2}\\
        \overline{\kappa} = \kappa\sqrt{1-p_i}\\
        \overline{\tau} =  \sqrt{\dfrac{\Gamma(\overline{\nu})}{\Gamma(\overline{\nu}+\sfrac{1}{2})\sqrt{4\pi}\overline{\kappa}^{2\overline{\nu}}\sigma^2}}\\
        \cdot\texttt{Qalpha1}(\overline{\tau}, \overline{\kappa})
    \end{cases}
\end{align}
\subsection{Case $\alpha\in(1,2)$}

\begin{align}
   \varrho_0(h)= k\dfrac{c_\alpha}{c_{\lfloor\alpha\rfloor}}\varrho(h;\kappa, \lfloor\alpha\rfloor-\sfrac{1}{2}, \sigma)\implies\mathbf{Q}_0 = \dfrac{1}{kc_\alpha/c_{\lfloor\alpha\rfloor}} \begin{cases}
        \overline{\nu} = \sfrac{1}{2}\\
        \overline{\kappa} = \kappa\\
        \overline{\tau} =  \sqrt{\dfrac{\Gamma(\overline{\nu})}{\Gamma(\overline{\nu}+\sfrac{1}{2})\sqrt{4\pi}\overline{\kappa}^{2\overline{\nu}}\sigma^2}}\\
        \cdot\texttt{Qalpha1}(\overline{\tau}, \overline{\kappa})
    \end{cases}
\end{align}
\begin{align}
    \varrho_i(h) =  \varrho_i(h) \implies \mathbf{Q}_i= \dfrac{1}{r_i\sigma^2}\left(\texttt{matern.p.precision}(\kappa,p_i,\alpha, \ell_e) - f(\texttt{matern.p.joint}(\kappa,p_i,\alpha))\right)
\end{align}

\bibliographystyle{plain}
\bibliography{references} 
\end{document}
